\chapter{Organizing larger applications}

In an application with only a few routes, it may still be convenient to keep every model, query, and handler in a single \code{main.rw}. In a wiki, CMS, or line-of-business system, however, \emph{structure} soon matters more than the syntax of individual language features. RWLang provides two complementary tools: \code{mod} organizes source files, while \code{object} defines the namespace of a domain API.

\section{Three separate questions}

Treat the structure of a larger application as three distinct layers:

\begin{center}
\begin{tabularx}{\textwidth}{@{}l Y@{}}
\toprule
\textbf{Tool} & \textbf{Responsibility} \\
\midrule
\code{mod} & In which file or source module does the code live? \\
\code{object} & Which business concept owns the model, query, page, or action? \\
\code{route} & What HTTP surface and what auth/cache/rate-limit policy expose it? \\
\bottomrule
\end{tabularx}
\end{center}

This separation is deliberate. A domain object does not automatically receive a database, current user, filesystem, or network access, and a module does not imply HTTP policy.

\section{\code{mod} is not textual include}

A small multi-file project might look like:

\begin{lstlisting}
myapp/
|-- main.rw
|-- models.rw
|-- queries.rw
|-- components.rw
|-- pages.rw
|-- actions.rw
|-- migrations/
`-- public/
\end{lstlisting}

The entry point:

\begin{lstlisting}[language=RWLang]
mod models;
mod queries;
mod components;
mod pages;
mod actions;
\end{lstlisting}

\code{mod} is not C-style textual inclusion and it does not inject declarations into a global namespace. The compiler loads, validates, and compiles the explicit source graph into one program. A module owns a namespace: declarations from \code{queries.rw}, for example, are referenced as \code{queries::name} across module boundaries. Module paths are application-root relative, so nested source files must be loaded explicitly with paths such as \code{mod article::admin;}.

\section{When should you switch to feature-based structure?}

Splitting by technical type -- \code{models.rw}, \code{queries.rw}, \code{pages.rw} -- is a good first step, but in a larger application it can mean opening five or six files for one feature change. At that point, organize by business area instead:

\begin{lstlisting}
main.rw
article.rw
article/
|-- admin.rw
`-- components.rw
account.rw
account/
`-- profile.rw
invoice.rw
invoice/
`-- admin.rw
routes.rw
migrations/
public/
\end{lstlisting}

Then \code{main.rw} becomes mostly application composition:

\begin{lstlisting}[language=RWLang]
mod article;
mod article::admin;
mod article::components;
mod account;
mod account::profile;
mod invoice;
mod invoice::admin;
mod routes;
\end{lstlisting}

The goal is not to maximize file count. The goal is for one business change to remain mostly inside one feature directory.

\section{Domain object: a business namespace, not an OOP class}

\code{object} groups related business names inside a domain namespace. For example:

\needspace{30\baselineskip}
\begin{lstlisting}[language=RWLang]
object Article {
    model {
        id: Int
        slug: Slug
        title: String
        body: String
        authorUsername: String
    }

    query fn bySlug(db: Db, slug: Slug)
        -> Result<Article, DbError> sql {
        SELECT id, slug, title, body, authorUsername
        FROM articles
        WHERE slug = :slug
    }

    page fn show(ctx: PageContext, db: Db, slug: Slug)
        -> Result<Html, PageError> {
        let article = Article.bySlug(db, slug)?;
        return Ok(html {
            <article>
                <h1>{{ article.title }}</h1>
                @markdown(article.body)
            </article>
        });
    }
}
\end{lstlisting}

Usage becomes expressive:

\begin{lstlisting}[language=RWLang]
let article = Article.bySlug(db, slug)?;
\end{lstlisting}

Short domain names become natural too, for example \path{Invoice.approve}, \path{Invoice.cancel}, and \path{UserAccount.disable}.

\subsection*{What may an object contain?}

A V1 domain object can group operations around one model. The following is only a structural sketch, not standalone compilable RWLang source:

\begin{lstlisting}
object Article {
    model { ... }
    query fn bySlug(...) ...
    page fn show(...) ...
    action fn publish(...) ...
    component fn Card(...) ...
    layout fn EditorLayout(...) ...
}
\end{lstlisting}

Exactly one model block belongs to an object. An object is not a classic object-oriented class: there is no inheritance, mutable \code{self}, constructor magic, virtual dispatch, or reflection.

\section{Dependencies remain explicit}

The domain namespace does not hide capabilities. This is good:

\needspace{11\baselineskip}
\begin{lstlisting}[language=RWLang]
query fn byId(db: Db, id: Int)
    -> Result<Article, DbError> sql {
    SELECT id, slug, title, body, authorUsername
    FROM articles
    WHERE id = :id
}
\end{lstlisting}

An action receives the capabilities it needs just as explicitly. The following is intentionally abbreviated structural code; the concrete typed state-changing query is shown in the CRUD and business-audit sections:

\needspace{12\baselineskip}
\begin{lstlisting}
action fn publish(ctx: ActionContext, db: Db, id: Int)
    -> Result<Redirect, PageError> {
    let article = Article.byId(db, id)?;
    authorize article owner authorUsername or role Publisher;
    // ... state transition inside a transaction ...
    return Ok(redirect("/admin/articles"));
}
\end{lstlisting}

The handler signature still reveals database access. The same principle applies to AppFs, auth, and later network egress: a domain object must not become a hidden service locator.

\section{Why should routes remain top-level?}

Routes deliberately do not disappear entirely into objects:

\begin{lstlisting}[language=RWLang]
route articleShow GET "/article/:slug<Slug>"
    cache public ttl 60
    => article::Article.show;

route articlePublish POST "/admin/article/:id<Int>/publish"
    auth role Publisher
    invalidate cache articleShow
    => article::Article.publish;
\end{lstlisting}

A route registry or code review can then see in one place:

\begin{itemize}
  \item URL and HTTP method;
  \item authentication/authorization requirement;
  \item cache policy and invalidation;
  \item rate limit or resource profile;
  \item the global route name used in logs and audit.
\end{itemize}

A domain member name may be short and locally meaningful, such as \code{Article.publish} inside its own module; from a separate \code{routes.rw}, use \code{article::Article.publish}. The route name should remain globally descriptive, such as \code{articlePublish}.

\section{Recommended naming convention}

\begin{center}
\begin{tabularx}{\textwidth}{@{}l l Y@{}}
\toprule
\textbf{Element} & \textbf{Form} & \textbf{Example} \\
\midrule
object & PascalCase business noun & \code{Article}, \code{Invoice}, \code{UserAccount} \\
member & camelCase, short inside domain & \code{bySlug}, \code{publish}, \code{approve} \\
route & globally auditable name & \code{articleShow}, \code{invoiceApprove} \\
module & feature or technical unit & \code{article}, \code{account}, \code{routes} \\
\bottomrule
\end{tabularx}
\end{center}

\section{One possible CMS structure}

A medium-size CMS might use:

\needspace{18\baselineskip}
\begin{lstlisting}
main.rw
routes.rw
article.rw            # Article object + base queries
article/
|-- admin.rw          # edit/publish
`-- components.rw     # Card, lists
media.rw
media/
`-- admin.rw
account.rw
migrations/
public/
`-- app.css
\end{lstlisting}

Public and admin HTTP policy remain in the route layer. The article feature keeps the article business model and its operations together.

\section{What not to do}

\subsection*{Do not create one file per declaration}

Overly small modules create more navigation than structure. A fifty-line feature does not need eight files.

\subsection*{Do not build a hidden ``service'' layer}

If an operation uses DB, AppFs, or auth context, keep that visible in the typed signature. One of RWLang's strengths is precisely that capability boundaries are explicit.

\subsection*{Do not duplicate route policy inside handlers}

Route-level auth, cache, rate limit, and resource policy are declarative contracts. Domain-specific object authorization naturally remains in the handler because it depends on the business record.

\subsection*{Do not use free-form String for business state}

If a domain has \code{Draft}, \code{Review}, \code{Published}, or \code{Cancelled} states, use an enum. Project structure helps most when the domain contract is typed too.

\section{Refactoring path: from one file to a larger application}

You do not need to design the final directory hierarchy on day one. A good evolutionary path is:

\begin{enumerate}
  \item Start with one \code{main.rw} while the application is small.
  \item Separate model/query/page/action groups when the file becomes hard to navigate.
  \item If one feature is scattered across many technical files, switch to feature-based modules.
  \item If too many global business names appear, group them under \code{object} namespaces.
\end{enumerate}
